\section{Problem and Experimental Setting}
\label{sec:setting}

\paragraph{Aliased T-maze.}
An episode starts in a clue cell whose observation reveals a hidden binary clue $c\in\{A,B\}$. The agent then walks a corridor of length $L$, sampled uniformly from $\{\DiagCfgCorridorMin,\dots,\DiagCfgCorridorMax\}$. Corridor observations are drawn from \DiagCfgDistractors{} distractor symbols independently of $c$. The corridor ends at a junction whose observation is identical for both clues. The actions are \textsc{forward}, \textsc{left}, and \textsc{right}. In the corridor, \textsc{left} and \textsc{right} leave the agent in place; at the junction, \textsc{forward} leaves the agent in place, and \textsc{left} or \textsc{right} ends the episode in an arm. In the symmetric variant used throughout, the arm that matches the clue gives a goal observation and reward $+1$, and the other arm gives reward $-1$. Episodes time out after $L+2+\DiagCfgTimeoutSlack$ steps. Two histories that differ only in the clue therefore share the current observation at the junction, yet the consequence of \textsc{left} differs between them. Under any clue-independent distribution over action sequences, the action-marginal distribution of future observations and rewards is identical for both clues. Our tests confirm this by exact enumeration for horizons up to three, starting at the junction and one step before it, and also confirm that the action-conditional futures differ. This is a direct instance of the PSR observation that the relevant predictions are action-conditional \citep{littman2001psr,singh2004psr}, and we do not present it as a new result.

\paragraph{Information access.}
The simulator's complete state is private: clue, corridor length, position, and random-number-generator state. It is available only to data collectors (for snapshot and restore) and to the evaluator (for labels). Actors, encoders, prediction heads, and the planner receive only the observable history $h_t=(o_0,a_0,r_1,\dots,a_{t-1},r_t,o_t)$. Ground-truth labels are returned in a separate structure that training code never receives. Unit tests check the relevant function signatures, the whitelisted record fields, that serialized datasets contain no hidden fields, and that the encoder output at time $t$ is unaffected by later inputs.

\paragraph{Transition budget and the restore privilege.}
Every environment step, whether on the main trajectory, in a branch rollout, or during evaluation, is charged to a ledger \emph{before} it executes. A step that would exceed the cap is not executed. Data collection spends the budget exactly, so arms with equal caps also have equal \emph{actual} transition spend. Snapshots and restores are not transitions, but they are a simulator privilege; the ledger counts them separately, and the evaluation ledger forbids them. An equal transition budget does not imply equal compute, and we report arithmetic cost (multiply-accumulates, MAC) and CPU time separately.

\paragraph{Branching collection.}
Main trajectories follow a history-based behaviour policy. It is uniform at the junction and presses \textsc{forward} with probability \DiagCfgForwardProb{} elsewhere. At each main step, with probability \DiagCfgBranchProb{}, the collector snapshots the simulator state. For each of the three first actions it then rolls out up to $k=\DiagCfgHorizon$ steps with a uniform continuation policy, restoring the snapshot between rollouts and after the last one. Random streams are separated so that branching never changes the main trajectory; a test checks this property. A branch set is started only if the remaining budget can cover it.

\paragraph{Readout protocol.}
Representations are compared through one fixed downstream path. The encoder is frozen. Its features are z-scored per dimension, with statistics fitted only on the readout training data. A fresh action-conditional reward head, a latent rollout with hidden size \DiagCfgHidden, is trained with squared reward error at learning rate \DiagCfgHeadLR. The head initialization, data, and optimizer are identical across encoders. We keep the checkpoints after \DiagCfgHeadEpochs{} epochs of a single run and select the one with the lowest development reward error; the test split is never used for selection. A model-predictive planner enumerates all action sequences of length \DiagCfgPlanHorizon{} and scores them by discounted predicted reward ($\gamma=\DiagCfgGamma$).

\paragraph{Metrics.}
Each test episode starts with a scripted, observation-only prefix that presses \textsc{forward} until the junction observation appears. \emph{Forced-commit accuracy}, the primary metric, is the fraction of test episodes in which the planner, restricted to first actions \textsc{left} and \textsc{right} at the first junction arrival, prefers the optimal arm. It measures one aliased decision and is not a planning success rate. The \emph{full-action} planner then acts with all three actions until the episode ends. From it we report success, mean return, and timeout rate, together with the number of exact forced-commit ties. The evaluator assigns clues by alternating episode index, so every constant decision rule scores exactly 0.5. A logistic probe of the clue on the junction features is fitted on a separate calibration split. We use it only as a diagnostic: failure of a linear probe does not show that the features contain no clue information.
